\scanpage{89}
\mspage{\handwritten{542} \unsure{\deleted{\unsure{5|3|B|6}}|\deleted{\unsure{X|3}}\unsure{5|B|6}}}

\noindent § 7. Ensembles constructibles dans les schémas\dotfill

7.1. Morphismes de présentation finie et ensembles constructibles\dotfill

7.2. Ensembles pro-constructibles et ind-constructibles.

7.3. Application aux morphismes générisants\dotfill

\noindent § 8. Applications rationnelles \dotfill

8.1. Applications rationnelles et fonctions rationnelles

8.2. Domaine de définition d'une application rationnelle

8.3. Faisceau des fonctions rationnelles\dotfill

8.4. Faisceaux de torsion et faisceaux sans torsion\dotfill

8.5. Critères de séparation d'un schéma intègre\dotfill

\noindent § 9. Foncteurs représentables élémentaires dans la théorie des schémas : Schémas relativement affines , fibrés vectoriels , fibrés projectifs et grassmanniennes..

9.1. Morphismes affines et spectres d'Algèbres quasi-cohérentes\dotfill

9.2. Faisceaux quasi-cohérents sur un $S$-schéma affine sur $S$\dotfill

9.3. Application aux changelents de base pour les images directes de Modules quasi-cohérents\dotfill

9.4. Fibré vectoriel associé à un faisceau de modules\dotfill

9.5. Schémas en ensembles algébriques\dotfill

9.6. Schémas en groupes linéaires\dotfill

9.7. Grassmanniennes d'un Module\dotfill

9.8. Morphismes de Plücker et de Segre\dotfill

9.9. F\deleted{\unsure{u|o|i}}ibrés en drapeaux\dotfill

9.10.Fibrés en variétés de Stiefel\dotfill

\noindent § 10. Schémas formels\dotfill

10.1. Schémas formels affines\dotfill

10.2. Morphismes de schémas form\unsure{\deleted{l}e|\deleted{e}l}\unsure{\deleted{e}l|\deleted{s}l|\deleted{l}e|\deleted{l}s}s affines\dotfill

10.3. Idéaux de définition d'un schéma formel affine\dotfill

10.4. Schémas formels et morphismes de schémas formels\dotfill

10.5. Idéaux de définition des schéma\unsure{\deleted{q}s|\deleted{s}q} formels\dotfill

10.6. Schémas formels comme limite inductives de schémas

10.7. Produit de schémas formels\dotfill

10.8. Complété formel d'un schéma le long d'un sous-schéma\dotfill
% archive sheet number: 88
